In the first stage, the router inspects the front flit of each input virtual-channel FIFO. The front flit can be checked without being immediately removed from the buffer. If the front flit is a head flit and the corresponding input VC is not already locked by a previous packet, the router decodes the destination information from the header and performs XY routing to determine the requested output port.

After the output port is determined, the router tries to allocate an available output virtual channel at that port. If the allocation succeeds, the selected output port and output VC are stored in the first-stage pipeline registers. At the same time, an input-VC lock is established. This lock records the reserved output port and output VC, so that the body and tail flits of the same packet can follow the same path without repeating routing computation and virtual-channel allocation.
